\documentclass[10pt,a4paper]{amsart}
\usepackage[margin=19mm]{geometry}
\usepackage{amsmath,amssymb,mathtools}
\usepackage[scr=rsfs,cal=euler]{mathalfa}
\usepackage{xfrac}
\usepackage[unicode,hidelinks,bookmarks=false]{hyperref}
\newcommand{\ie}{i.e.\ }
\newcommand{\cf}{cf.\ }
\newcommand{\resp}{respectively\ }
\newcommand{\Subsection}{\subsection}
\providecommand{\nobreakdash}{\nobreak\hbox{-}}
\setlength{\emergencystretch}{2em}
\multlinegap0pt
% ---------------------------------------------------------------------------
% Editorial AMS wrapper prelude. The theorem environments and the mathematical macros below are
% transcribed from the paper's own preamble (cached-originals/source0.tex lines 1--38); the AMS amsart
% class, the named format packages of the pinned TeX Live runtime and this editorial formatting are not
% part of the author's text. No mathematics is changed.
% ---------------------------------------------------------------------------
\newtheorem{theorem}{Theorem}[section]
\newtheorem{proposition}[theorem]{Proposition}
\newtheorem{lemma}[theorem]{Lemma}
\newtheorem{corollary}[theorem]{Corollary}
\theoremstyle{remark}
\newtheorem{remark}[theorem]{Remark}
\numberwithin{equation}{section}
\newcommand{\Z}{\mathbb Z}
\newcommand{\E}{\mathbb E}
\newcommand{\PP}{\mathbb P}
\newcommand{\HH}{\mathcal H}
\newcommand{\LL}{\mathcal L}
\newcommand{\one}{\mathbf 1}
\DeclareMathOperator{\Bin}{Bin}
\DeclareMathOperator{\rk}{r}
\DeclareMathOperator{\diver}{div}

\begin{document}
\title{Paths maximize the expected range of graph-indexed random walks}
\author{Yinfeng Zhu}
\date{Source: arXiv:2609.19728v1; distributed under CC BY 4.0}
\maketitle
\noindent\emph{Source, version and licence.} \emph{Paths maximize the expected range of graph-indexed random walks}, by Yinfeng Zhu. The source of this editable unit is the e-print of \textbf{version v1} of arXiv:2609.19728, https://arxiv.org/abs/2609.19728v1, whose material is distributed under the Creative Commons Attribution 4.0 International licence, http://creativecommons.org/licenses/by/4.0/, as recorded in the arXiv licence metadata for this paper and shown on that abstract page. The whole of the body below is version v1, the newest version of the paper. Full rights notice: licenses/arxiv-2609.19728-CC-BY-4.0.txt.\par\smallskip
\noindent\textit{Editorial scope.} Counted unit: original Theorem 1.1 of the article, the statement carried by the source environment \texttt{theorem} with source label \texttt{thm:bhm}, cached-originals/source0.tex lines 99--104, printed as \textbf{Theorem 1.1} exactly as in the article: the statement that among all graphs on a given number of vertices the path maximizes the expected range of the graph-indexed random walk. It is delivered together with the authors' own complete proof as the single primary proof body, source0.tex lines 931--975 (45 lines): the \texttt{proof} environment that belongs to the statement and establishes it in full. No proof marker is added and no mathematics is written, repaired or re-derived; the authors' notation and the printed slips of the source are kept literal. Necessary complete context is retained uncounted, in source order: lines 269--271, the source \texttt{equation} with source label \texttt{eq:decimation} of the article that the retained passages refer to; lines 900--908, the source \texttt{proposition} with source label \texttt{prop:single-class} of the article that the retained passages refer to. The article's own numbering is reproduced by explicit counter settings: the theorem-like counter of the article is set before each retained passage, so the counted statement prints with the number the article gives it. No \texttt{\textbackslash cite} command occurs in any retained passage: the closure of the retained passages over references and citations was computed, it contains no citation at all, so this unit delivers no bibliography entry and the article's own bibliography data is neither spliced into the bundled source nor redistributed. The source's own class is \texttt{amsart}; the e-print ships no class or style file, so no publisher class is shipped, used or redistributed, and the wrapper is the AMS \texttt{amsart} class with the article's own macros and theorem declarations transcribed from its preamble. The retained passages contain no figure environment and no graphics-inclusion command, so no artwork is redistributed. Every declared body of this unit is an exact, unmodified span of source0.tex: no body is a bounded passage comparison, \texttt{omit} was not used anywhere, and consequently no fidelity receipt is asserted in delivery-evidence.json. The complete concatenated original source is bundled as sources/210804/source0.tex. Omitted from this editable unit and therefore absent from it are source0.tex lines 1--98; 105--268; 272--899; 909--930; 976--1331. The AMS \texttt{amsart} wrapper, the named format packages of the pinned TeX Live runtime and this editorial source, licence and scope paragraph are editorial formatting only.\par\smallskip
\setcounter{section}{1}
\setcounter{equation}{0}
\setcounter{theorem}{0}
\begin{theorem}[BHM expectation conjecture]\label{thm:bhm}
For every finite connected bipartite simple graph $G$ on $n$ vertices,
\begin{equation}\label{eq:bhm}
 \widehat h(G)\le\widehat h(P_n).
\end{equation}
\end{theorem}

\setcounter{section}{2}
\setcounter{equation}{3}
\setcounter{theorem}{3}
\begin{equation}\label{eq:decimation}
 \widehat h(G)=1+\E\widehat h(Q_B)+\E\widehat h(Q_W).
\end{equation}

\setcounter{section}{4}
\setcounter{equation}{11}
\setcounter{theorem}{3}
\begin{proposition}[Single-class comparison]\label{prop:single-class}
For either bipartition class, with $d$ one less than its size and $K$
the number of components of its full zero-edge subgraph,
\begin{equation}\label{eq:single-class}
 \E b_{K-1}\le\E b_{Y_{d,1/2}}.
\end{equation}
The inequality is strict if the corresponding auxiliary graph contains
a cycle.
\end{proposition}

\setcounter{section}{4}
\setcounter{equation}{12}
\setcounter{theorem}{4}
\begin{proof}[Proof of Theorem~\ref{thm:bhm}]
We proceed by induction on $n=|V(G)|$. If $n=1$, both ranges are zero.
Suppose
$n\ge2$ and that the theorem holds for all smaller connected bipartite
graphs. Let $|B|=a$ and $|W|=b$, so $a+b=n$. Put
\[
 t_j=\E b_{Y_{j,1/2}}\qquad(j\ge0).
\]
Each quotient in~\eqref{eq:decimation} has fewer than $n$ vertices, so
the induction hypothesis applies to every realized quotient. Applying
Proposition~\ref{prop:single-class} then gives
\begin{equation}\label{eq:partition-bound}
 \widehat h(G)
 \le1+\E b_{K_B-1}+\E b_{K_W-1}
 \le1+t_{a-1}+t_{b-1}.
\end{equation}

To compare the two bipartition sizes, we use the concavity of $(t_j)$.
Couple $Y_{j+1,1/2}=Y_{j,1/2}+\xi$, where the additional fair Bernoulli
variable $\xi$ is independent. This coupling gives
\[
 t_{j+1}-t_j
   =\frac12\E\bigl[b_{Y_{j,1/2}+1}-b_{Y_{j,1/2}}\bigr].
\]
The increments of $b$ are nonincreasing, and the coupled binomial
variables are nondecreasing in $j$. Hence these increments of $t$ are
nonincreasing. Balancing two nonnegative indices with sum $n-2$ gives
\begin{equation}\label{eq:balance}
 t_{a-1}+t_{b-1}
 \le t_{\lfloor(n-2)/2\rfloor}+t_{\lceil(n-2)/2\rceil}.
\end{equation}

To identify the balanced bound, apply the exact
decomposition~\eqref{eq:decimation} to $P_n$.
Its two auxiliary graphs are paths, including the possible one-vertex
path, and every auxiliary edge has $\lambda_e=2$. Their nonzero-edge
counts are binomial with parameter $1/2$, and each zero-edge quotient
is again a path. Thus
\begin{equation}\label{eq:path-decimation}
 b_{n-1}=1+t_{\lfloor(n-2)/2\rfloor}
           +t_{\lceil(n-2)/2\rceil}.
\end{equation}
Combining~\eqref{eq:partition-bound}--\eqref{eq:path-decimation} proves
$\widehat h(G)\le b_{n-1}=\widehat h(P_n)$ and completes the induction.
\end{proof}

\end{document}
